% !TeX root = ../thesis.tex
\section{Introduction}
\label{sec:introduction}

\subsection{Background and Motivation}
\label{sec:introduction-background}

Text-to-SQL systems allow users to interact with a database through natural language. 
A model translates users' questions into SQL statements, 
which can then be executed to retrieve the relevant data. 
Research using Large Language Models (LLMs) has developed several approaches to this task, 
including prompting and fine-tuning \citep{shi2025survey}.

In an agentic system, answering a question and returning the final result 
can involve multiple interactions with the underlying database system. 
The agent may generate one or more queries, 
possibly encounter errors, and retry, before finally providing an answer.
A correct final result can therefore follow an unsuccessful first attempt. 
Similarly, a failed initial SQL attempt may be followed by several additional retries 
before the last query executes successfully and the result is returned to the user.
Hence, evaluating only the final answer gives an 
indication of success or failure, but leaves the preceding attempts unexplained.

This distinction matters when examining how a deployed system behaves and handles errors.
When considering only final accuracy, both a correct result reached on the first attempt 
and a correct result reached after several failures are treated equally.
There is therefore value in saved SQL attempts and database feedback 
because they can provide evidence for examining these behaviours.

\subsection{Research Problem}
\label{sec:introduction-problem}

Models' behaviour beyond final accuracy 
is already the subject of ongoing research.
For example, \citet{yan2026bottleneck} use trajectory analysis 
to identify gaps in Text-to-SQL systems and develop a framework 
for collecting required information before query generation.

The literature reviewed in
\hyperref[sec:literature-review]{Chapter~\ref*{sec:literature-review}}
also covers aspects such as SQL correction, intermediate progress, 
conversational evaluation and related topics.

At the centre of this study is the question of 
how recorded SQL attempts can be used 
to explain performance of a deployed Text-to-SQL workflow.

In particular, the study connects these attempts 
with database feedback and reference-result comparisons. 
This allows the final outcome to be examined alongside SQL query revisions 
and other events leading to a successful answer or a failure.

The comparison between two selected models was chosen as the main investigation method.
This makes it possible to identify whether models answering the same questions 
can provide the same answers and reach a similar result through different trajectories.
Additionally, it enables examination of whether outcomes and classes remain stable across repeated runs.
Therefore, the purpose is to understand what can be inferred from the traces, 
especially with regard to the system's behaviour and limitations.

\subsection{Research Aim and Questions}
\label{sec:research-questions}

The overall aim of this study, therefore, is to determine how trajectory-aware evaluation changes
the interpretation of final result accuracy in an agentic Text-to-SQL system.

Specifically, this work investigates the following primary research question:
\begin{quote}
To what extent does trajectory-aware evaluation change the interpretation
of final result accuracy by distinguishing direct success, recovery and
unresolved failure, and how do Nova Lite and Haiku 4.5 differ in their
recovery, clarification, error-response and failure behaviour?
\end{quote}

To address this question, the analysis is structured around five supporting sub-questions:
\begin{itemize}
  \item \textbf{RQ1:} What proportion of successful final outcomes were
  not achieved through clean direct success?
  \item \textbf{RQ2:} How much final task success involved recovery from
  an initially unsuccessful SQL attempt?
  \item \textbf{RQ3:} How do the agents respond to database execution
  errors, and which observed patterns are associated with recovery
  or repeated failure?
  \item \textbf{RQ4:} How do correctness, trajectory classes and
  SQL-attempt counts vary across the benchmark's difficulty groups?
  \item \textbf{RQ5:} When the models reach the same final-outcome
  category on the same question, how often do their trajectory classes differ?
\end{itemize}

\subsection{Study Scope and Contributions}
\label{sec:introduction-scope}

The study used a cloud-based Text-to-SQL system to conduct the experiments 
and gather the data required for the analysis.
At its core, the system consists of Amazon Bedrock, AWS Lambda and Amazon Athena.
The models selected for the experiments were Amazon Nova Lite and Claude Haiku 4.5.
Each completed two runs of a project-specific benchmark based on the Chinook dataset.

The benchmark included standalone one-off and ambiguous questions,
two-turn interactions testing the continuation of an earlier conversation,
and safety-related requests designed to elicit inappropriate responses and test the models' refusal behaviour.
All query results and execution traces were saved, 
and used to assess correctness and to reconstruct SQL trajectories. 

Based on this investigation, this study makes three primary contributions:
\begin{itemize}
  \item An evaluation methodology connecting final correctness with
  SQL trajectories, error responses and recovery.
  \item Software that reconstructs and evaluates the recorded attempts
  using saved evidence and reference results.
  \item An empirical analysis of differences across models, difficulty
  groups and repeated runs, supported by selected examples.
\end{itemize}

The scope of the research and the evaluation performed is restricted to 
observable execution actions and results within the selected database, benchmark, and runs.
The models' hidden reasoning, the quality of their natural-language answers 
and their usefulness to the users were not examined and remain out of scope of this work.
However, the developed software could support further evaluations if 
the required evidence formats and comparison contracts are preserved.

All source code, including the LaTeX source of this thesis, is available at
\url{https://github.com/MateuszKikmunter/msc-research}.

\subsection{Thesis Structure}
\label{sec:introduction-structure}

The structure of this thesis is as follows.
\hyperref[sec:literature-review]{Chapter~\ref*{sec:literature-review}}
reviews the literature on Text-to-SQL correctness, repair, trajectory evaluation,
and interaction.
\hyperref[sec:architecture]{Chapter~\ref*{sec:architecture}}
presents the deployed system architecture and implementation, followed by
\hyperref[sec:methodology]{Chapter~\ref*{sec:methodology}},
which details the experimental design and evaluation methodology.
\hyperref[sec:results]{Chapter~\ref*{sec:results}}
reports the experimental results, while
\hyperref[sec:discussion]{Chapter~\ref*{sec:discussion}}
discusses their interpretation, limitations, and implications. Finally,
\hyperref[sec:conclusion]{Chapter~\ref*{sec:conclusion}}
concludes the thesis with a summary of key findings.
